\subsection{Подготовка изображений и разбиение}
\textbf{Первый проход:} откройте одно изображение, найдите его метку и форму тензора, затем проверьте размер выхода модели. До длительного обучения оцените память и выполните короткий пробный проход. Он проверяет запуск, но ещё не показывает качество обучения.

\subsubsection*{Шаг 1. Изучите набор и проверьте ресурсы}
Выведите пять классов, число изображений каждого класса и несколько исходных примеров. Проверьте читаемость файлов, размеры, цветовые каналы и дубликаты. Не задавайте число изображений по чужому описанию: оно зависит от версии выгрузки. Выберите устройство и запишите его в отчёте.
\begin{lstlisting}
import torch
from lab3_utils import prepare_flowers, seed_everything
from lab3_utils import original_cnn, transfer_resnet18, compact_example
from lab3_utils import fit_flowers, plot_learning, parameter_counts
seed_everything(0)
device = "cuda" if torch.cuda.is_available() else "cpu"
print(device)
\end{lstlisting}

\subsubsection*{Шаг 2. Зафиксируйте одну валидацию}
Исходный ноутбук выделяет 1000 изображений. В исправленном варианте это ориентир размера валидации; точные копии файлов объединяются по SHA-256 и попадают в одну часть. Поэтому фактическое число проверочных изображений может отличаться от 1000. \file{split.csv} сохраняет итоговое распределение. Групповое разделение не гарантирует одинаковые пропорции классов: выведите их и проверьте наличие всех пяти классов в обеих частях.
\begin{lstlisting}
data = prepare_flowers("data", validation_target=1000,
                       seed=0, image_size=224, augment=False)
print(data["class_to_idx"])
print(data["inventory"].groupby(["split", "target"]).size())
print(data["config"])
\end{lstlisting}

Хеши находят побайтовые копии. Снимки одной сцены после повторного сжатия, обрезки или изменения размера могут иметь разные хеши. Если такие серии известны, их также следует объединить до разделения. Совпадающие файлы с разными метками требуют проверки разметки. Не исправляйте конфликты только ради улучшения метрики.

Обучающий и проверочный наборы создаются как отдельные \file{ImageFolder} с одним списком файлов и разными подмножествами. При дополнительном опыте со случайным отражением оно применяется только к обучению. Замена \file{transform} у общего объекта под двумя \file{Subset} изменила бы обе части одновременно.

\begin{keybox}{Контрольная точка 1}
Пакет имеет форму \((B,3,224,224)\), метки --- целочисленный вектор длины \(B\) со значениями от 0 до 4. Известно отображение индексов в имена. Обучение и валидация не пересекаются по файлам и точным дубликатам. Все три архитектуры используют одно сохранённое разбиение.
\end{keybox}

\subsection{Обучение и исследование архитектур}
\subsubsection*{Шаг 3. Разберите и обучите заданную CNN}
Выпишите размеры тензоров после свёрток и pooling, проверьте вход первого Linear и число параметров. Создайте новую модель, затем оптимизатор. Функция ниже использует Adam с явно переданным шагом, десять эпох и пакет 32. Её код доступен в \file{lab3\_utils.py}; найдите места сброса метрик, переключения режима и выбора сохранённых весов.
\begin{lstlisting}
seed_everything(0)
model = original_cnn()
print(parameter_counts(model))
result_original = fit_flowers(
    model, data, "runs/lab3_original_01", epochs=10,
    batch_size=32, learning_rate=1e-3, seed=0, device=device
)
\end{lstlisting}

Перед полным запуском проверьте один пакет. Если памяти недостаточно, уменьшите пакет и заново зафиксируйте общий протокол для всех трёх моделей. Уменьшение разрешения у исходной сети потребует изменения первого Linear; это уже другая архитектура и должно быть указано. Не заносите ошибку памяти в таблицу качества как нулевую accuracy.

\subsubsection*{Шаг 4. Примените перенос обучения}
Выберите предобученную архитектуру, прочитайте её описание и определите слой классификатора. Пример ниже конкретизирует ResNet18. Порядок важен: сначала заморозка старых параметров, затем создание новой головы.
\begin{lstlisting}
from torchvision import models
from torch import nn
weights = models.ResNet18_Weights.IMAGENET1K_V1
model = models.resnet18(weights=weights)
for p in model.parameters():
    p.requires_grad = False
model.fc = nn.Linear(model.fc.in_features, 5)
model.frozen_backbone = True
for name, p in model.named_parameters():
    if p.requires_grad:
        print(name, p.shape)
\end{lstlisting}
Вывод должен содержать только \file{fc.weight} и \file{fc.bias}. Для головы \(512\to5\) обучается \(512\cdot5+5=2565\) параметров. При первом использовании веса могут загружаться из сети; \file{weights=None} создаёт случайную инициализацию и не является переносом обучения.

Фабрика \file{transfer\_resnet18()} выполняет эти действия и сохраняет обозначение весов в метаданных. Используйте новую модель и новый каталог опыта:
\begin{lstlisting}
seed_everything(0)
model = transfer_resnet18(pretrained=True)
result_transfer = fit_flowers(
    model, data, "runs/lab3_transfer_01", epochs=10,
    batch_size=32, learning_rate=1e-3, seed=0, device=device
)
\end{lstlisting}
В режиме фиксированного извлекателя проверьте, что параметры backbone и статистики BatchNorm не меняются после шага. Если исследуете fine-tuning, явно укажите размороженные слои, режим BatchNorm и новый оптимизатор. Не представляйте этот дополнительный опыт как тот же режим заморозки.

\subsubsection*{Шаг 5. Создайте собственную архитектуру}
Предложите изменение, связанное с проверяемой гипотезой: уменьшение головы, другая ширина блоков, изменение глубины, регуляризация или иной способ агрегирования карт. Нарисуйте последовательность слоёв, рассчитайте формы и число параметров. Выход обязан иметь пять логитов. Маленькая модель ниже иллюстрирует конструкцию с глобальным усреднением:
\begin{lstlisting}
student_model = nn.Sequential(
    nn.Conv2d(3, 16, 3, padding=1), nn.ReLU(), nn.MaxPool2d(2),
    nn.Conv2d(16, 32, 3, padding=1), nn.ReLU(), nn.MaxPool2d(2),
    nn.Conv2d(32, 64, 3, padding=1), nn.ReLU(),
    nn.AdaptiveAvgPool2d(1), nn.Flatten(), nn.Linear(64, 5)
)
\end{lstlisting}

\textbf{Самостоятельная часть:} внесите как минимум одно осмысленное структурное изменение относительно данного примера и объясните ожидаемый эффект. Простое копирование \file{compact\_example()} не считается разработкой своей архитектуры. Обучите свой вариант с тем же разбиением и основным бюджетом. Зафиксируйте seed \emph{до} создания слоёв, а не только перед вызовом функции обучения.

\subsection{Сравнение, визуализация и сохранение}
\subsubsection*{Шаг 6. Соберите сопоставимые результаты}
Основной протокол: RGB 224\(\times\)224, фиксированная нормализация, без случайной аугментации, seed 0, Adam с шагом \(10^{-3}\), пакет 32, десять эпох. Для каждой модели создаются новые загрузчики с тем же начальным seed. Выбирается эпоха с наибольшим macro-F1 валидации; при равенстве сохраняется более ранняя. Это одинаковое правило выбора, но не одинаковый объём предобучения.

\begin{table}[H]\centering\small
\caption{Что должно быть в итоговой таблице CNN}
\begin{tabularx}{\textwidth}{@{}p{0.30\textwidth}Y@{}}
\toprule
Группа полей & Содержание \\
\midrule
Модель & Имя, схема, число всех и обучаемых параметров; исходные веса \\
Обучение & Разбиение, seed, эпохи, пакет, оптимизатор, шаг, размер входа, аугментация \\
Результат & Лучшая эпоха, loss, accuracy и macro-F1 одной валидации \\
Ресурсы & Устройство, измеренное время с указанием границ измерения; при измерении --- пиковая память \\
Вывод & Какие классы и случаи улучшились или ухудшились; ограничения сравнения \\
\bottomrule
\end{tabularx}\end{table}

В модуле время включает обучение, валидацию и запись контрольных точек; оно не включает первоначальную загрузку предобученных весов и подготовку набора. Не сравнивайте его с измерением только одного прямого прохода. Одинаковое число эпох также не означает одинаковое число операций для разных сетей.

\subsubsection*{Шаг 7. Постройте графики и исследуйте прогнозы}
Для каждой модели покажите train/validation loss и macro-F1 по эпохам, а затем общую диаграмму результатов трёх архитектур. Сохраните матрицы ошибок и classification report по пяти классам. Используйте одинаковые оси и порядок меток, подпишите доли или проценты.
\begin{lstlisting}
fig = plot_learning(result_transfer)
fig.savefig("runs/lab3_transfer_01/curves.pdf", bbox_inches="tight")
\end{lstlisting}

Покажите девять заранее определённых изображений валидации либо все, если их меньше: истинный и предсказанный класс, восстановленные цвета. Дополнительно разберите несколько ошибок; не выбирайте только удачные картинки. Фрагмент демонстрирует корректный прогноз одного нормализованного изображения:
\begin{lstlisting}
from lab3_utils import denormalize
import matplotlib.pyplot as plt
index = int(data["validation_indices"][0])
image, label = data["validation_dataset"][index]
model = result_transfer["model"]
model.eval()
with torch.no_grad():
    logits = model(image.unsqueeze(0).to(device))
predicted = int(logits.argmax(1).item())
plt.imshow(denormalize(image).permute(1, 2, 0))
plt.title(f"True: {data['classes'][label]}; predicted: {data['classes'][predicted]}")
plt.axis("off")
\end{lstlisting}

\subsubsection*{Шаг 8. Сохраните воспроизводимый комплект}
Каждая архитектура получает свой каталог. Основные файлы:
\begin{itemize}
\item \file{split.csv} --- список изображений и разбиение;
\item \file{history.csv} --- история обучения по эпохам;
\item \file{best\_weights.pth} --- выбранные веса;
\item \file{manifest.json} --- параметры и результаты опыта;
\item \file{architecture.txt} --- текстовое описание модели.
\end{itemize}
Также сохраняются прогнозы, отчёт и матрица ошибок. Приложите исходный код собственной архитектуры. Для восстановления нужны совместимое определение модели, те же преобразования и тот же словарь классов, а не только файл весов.

Валидация используется для выбора эпохи и сравнения архитектур. Поэтому её лучший результат нельзя назвать независимой оценкой на новых фотографиях. Дополнительный финальный тест требует отдельного набора, не использованного при этих решениях.

\subsection{Контрольные вопросы и критерии готовности}
\begin{enumerate}
\item Чем CNN отличается от полносвязной сети по связям, параметрам и работе с пространством?
\item Что задают kernel size, stride, padding, входные и выходные каналы? Рассчитайте размер выхода одного слоя.
\item Как работают свёртка и pooling в исходной сети? Какие из них имеют обучаемые параметры?
\item Почему первый Linear исходной архитектуры требует так много памяти?
\item Что такое логиты? Почему перед CrossEntropyLoss не добавляется Softmax?
\item Зачем нужны одновременно eval и отключение градиентов?
\item Что такое предобучение, transfer learning и fine-tuning? Откуда взяты веса выбранной модели?
\item Как устроена выбранная предобученная архитектура? Для ResNet объясните короткий путь и сложение в блоке.
\item Что делает requires\_grad? Почему заморозка параметров не равна заморозке статистик BatchNorm?
\item Какое изменение вы внесли в собственную сеть и какую гипотезу проверяли?
\item Почему нельзя накапливать счётчики разных эпох или усреднять F1 пакетов как F1 всего набора?
\item Как корректно показать нормализованное изображение? Почему clip без обратной нормализации неверен?
\item Что можно и чего нельзя заключить из сравнения с предобученной сетью?
\end{enumerate}

\begin{keybox}{Перед сдачей работы №\,3}
Обучены и сопоставлены все три требуемые архитектуры; объяснены свёртка, pooling, перенос обучения и заморозка. Есть собственное изменение модели, общая таблица, графики и анализ ошибок. Преобразования и разбиение согласованы, метрики относятся к отдельным эпохам, а не ко всему предыдущему обучению.
\end{keybox}
